\section{相关工作}
\label{sec:related}

\subsection{雷达网络部署优化}

雷达网络部署优化问题最早由Kirubarajan等人提出，采用遗传算法求解多雷达布局问题。近年来，多目标优化方法被广泛应用于该领域。Chen等(2023)使用NSGA-II算法优化雷达网络布局，在覆盖率最大化的同时最小化部署成本。

然而，现有方法普遍忽视了边界效应问题和空地传播差异。边界处探测概率下降15\%-20\%的问题在文献中很少被系统性地解决。Kumar等(2022)指出，区域分解方法是解决边界效应的有效途径，但其算法复杂度较高，难以应用于大规模场景。

本文继承Han等(2025)的MOPSO-DT框架，该框架通过区域分解和坐标变换有效解决了边界效应问题。在此基础上，本文引入空地异构传播模型，使其更适用于空地协同电子对抗网络的场景。

\subsection{MOPSO算法}

多目标粒子群优化(MOPSO)由Coello等(2004)提出，通过引入外部档案存储非支配解，并使用拥挤距离维护解的多样性。该方法在多目标优化领域表现优异，被广泛应用于工程设计、路径规划等问题。

MOPSO-DT算法在标准MOPSO基础上增加了两个关键步骤：(1)区域分解：将复杂区域分解为凸多边形，便于坐标变换；(2)坐标变换：使用垂直交点法将归一化坐标映射到物理空间，消除边界效应。

本文采用Pareto支配关系确定个体最优和全局最优，并使用拥挤距离维护外部档案的多样性。实验表明，该方法在多目标优化问题上具有良好的收敛性和解多样性。

\subsection{空地协同传播模型}

自由空间传播模型假设电磁波在理想介质中传播，路径损耗与距离的平方成正比。然而，实际作战环境中，空-地(A2G)链路和地-地(G2G)链路的传播特性差异显著。

A2G链路具有接近自由空间的传播特性，路径损耗指数约为2.0-2.5，多径效应较小。这使得空中平台能够实现远距离探测和干扰。

G2G链路受地形遮挡和多径效应影响严重，路径损耗指数约为3.5-4.0。地面雷达的探测范围受限于视距和非视距传播条件。

国际电信联盟(ITU)在Recommendation ITU-R P.528中给出了A2G和G2G链路的路径损耗建议值。本文采用A2G链路的路径损耗指数lpha=2.0，G2G链路的路径损耗指数lpha=4.0，以反映两种链路的传播差异。